\documentclass[english, parskip=full]{scrartcl}
\usepackage[utf8]{inputenc}  % Kodierung der Datei
\usepackage[english]{babel}  % Sprache auf Deutsch 
\usepackage{color}
\usepackage{graphicx, gensymb}
\usepackage{amsfonts, cancel, amssymb}
\usepackage{amsmath, amsthm, array, esint}
\usepackage{booktabs, multirow}  % schönere Tabellen
\usepackage{caption}  % Anpassung der Captions
\usepackage{scrlayer-scrpage}    % Manipulation des Seitenstils
\pagestyle{scrheadings}  % Stil für die Seite setzen
\automark{section}  % Abschnittsnamen als Seitenbeschriftung 
\setheadsepline{.5pt}    % Kopzeile durch Linie abtrennen
\usepackage[hidelinks]{hyperref}  % Links und weitere PDF-Features
\usepackage{typearea, tikz, pgffor, verbatim}
\usetikzlibrary{calc, quotes}
\usepackage[cache=false, outputdir=build]{minted}
\usepackage{placeins} % provides \FloatBarrier

\definecolor{bg}{rgb}{0.9,0.9,0.9}
\newcommand\numberthis{\addtocounter{equation}{1}\tag{\theequation}}
\newcommand{\bra}{}
\newcommand{\kom}{}
\newcommand{\brak}{}
\newcommand{\braket}{}
\newcommand{\intx}{}
\newcommand{\intr}{}
\newcommand{\kla}{}
\newcommand{\klam}{}
\newcommand{\klammer}{}
\newcommand{\oper}{}
\newcommand{\tecxt}{}
\newcommand{\Bra}{}
\newcommand{\Ket}{}
\renewcommand{\i}{\text{i}}
\newcommand{\e}{\text{e}}
\newcommand*{\Fourier}{\textsc{Fourier}\ }
\newcommand*{\F}{\mathcal{F}}
\newcommand*{\sinc}{\text{sinc\,}}
\newcommand{\conv}{\mathop{\scalebox{3.5}{\raisebox{-0.38ex}{\makebox[0.5\width][c]{$\ast$}}}}\limits}

\newcommand*{\titel}{04 Bogoliubov-de Gennes Hamiltonians}
\newcommand*{\autor}{Joachim Schwardt}

\hypersetup{pdfauthor={\autor}, pdftitle={\titel}}

%\titlehead{titlehead}
\subject{Topological Insulators and Superconductors}
\title{\titel}
\author{\autor}
\date{\today}

\begin{document}

%\begin{titlepage}
\maketitle  % Titel setzen
%\tableofcontents  % Inhaltsverzeichnis setzen
%\end{titlepage}

\section*{Problem 4.1: Particle-Hole Symmetry}
We consider a particle-hole-symmetric Bloch Hamiltonian satisfying
\begin{align}
CH(\mathbf{p})C^\dagger &= -H(-\mathbf{p})^\ast,
\end{align}
where $C$ is the charge conjugation matrix.
Let the Hamiltonian take the form
\begin{align}
H(\mathbf{p}) &= \epsilon(\mathbf{p}) \tau_0 + d_\alpha (\mathbf{p}) \tau_\alpha
\end{align}
where $\tau_0$ is the identity matrix, and $\tau_\alpha$ are the Pauli-matrices.

\subsection*{a) $C=\tau_x$}
In this case, we have
\begin{align}
CH(\mathbf{p})C^\dagger &= \tau_x \kla\left( \epsilon(\mathbf{p}) \tau_0 + d_x (\mathbf{p}) \tau_x  + d_y (\mathbf{p}) \tau_y + d_z (\mathbf{p}) \tau_z \right) \tau_x = \\
&= \epsilon\tau_0 + d_x\tau_x - d_y\tau_y - d_z\tau_z \overset{!}{=} -H(-\mathbf{p})^\ast.
\end{align}
In explicit matrix form, we then find
\begin{align}
\begin{pmatrix}
\epsilon (\mathbf{p}) - d_z(\mathbf{p}) & d_x(\mathbf{p}) + \i d_y (\mathbf{p}) \\ 
d_x (\mathbf{p}) - \i d_y(\mathbf{p}) & \epsilon(\mathbf{p}) + d_z (\mathbf{p})
\end{pmatrix} &\overset{!}{=} -\begin{pmatrix}
\epsilon^\ast (-\mathbf{p}) + d_z^\ast(-\mathbf{p}) & d_x^\ast(-\mathbf{p}) + \i d_y^\ast(-\mathbf{p}) \\ 
d_x^\ast (-\mathbf{p})  - \i d_y^\ast(-\mathbf{p}) & \epsilon^\ast(-\mathbf{p}) - d_z^\ast (-\mathbf{p})
\end{pmatrix}.
\end{align}
We may assume that all coefficients are real, since we require $H$ to be hermitian, i.e. $H^\dagger = H$.
Then the most simple set of constraints is given by
\begin{align}
\epsilon (-\mathbf{p}) &= -\epsilon (\mathbf{p}),\\
d_x (-\mathbf{p}) &= -d_x (\mathbf{p}),\\
d_y (-\mathbf{p}) &= -d_y (\mathbf{p}),\\
d_z (-\mathbf{p}) &= d_z (\mathbf{p}).
\end{align}


\subsection*{b) $C=\tau_0$}
In this case, we have
\begin{align}
CH(\mathbf{p})C^\dagger &= \tau_0 \kla\left( \epsilon(\mathbf{p}) \tau_0 + d_\alpha (\mathbf{p}) \tau_\alpha \right) \tau_0 = \\
&= H(\mathbf{p}) \overset{!}{=} -H(-\mathbf{p})^\ast.
\end{align}
In explicit matrix form, we then find
\begin{align}
\begin{pmatrix}
\epsilon (\mathbf{p}) + d_z(\mathbf{p}) & d_x(\mathbf{p}) - \i d_y (\mathbf{p}) \\ 
d_x (\mathbf{p}) + \i d_y(\mathbf{p}) & \epsilon(\mathbf{p}) - d_z (\mathbf{p})
\end{pmatrix} &\overset{!}{=} -\begin{pmatrix}
\epsilon^\ast (-\mathbf{p}) + d_z^\ast(-\mathbf{p}) & d_x^\ast(-\mathbf{p}) + \i d_y^\ast(-\mathbf{p}) \\ 
d_x^\ast (-\mathbf{p})  - \i d_y^\ast(-\mathbf{p}) & \epsilon^\ast(-\mathbf{p}) - d_z^\ast (-\mathbf{p})
\end{pmatrix}.
\end{align}
We may assume that all coefficients are real, since we require $H$ to be hermitian, i.e. $H^\dagger = H$.
Then the most simple set of constraints is given by
\begin{align}
\epsilon (-\mathbf{p}) &= -\epsilon (\mathbf{p}),\\
d_x (-\mathbf{p}) &= -d_x (\mathbf{p}),\\
d_y (-\mathbf{p}) &= d_y (\mathbf{p}),\\
d_z (-\mathbf{p}) &= -d_z (\mathbf{p}).
\end{align}

%\begin{align}
%CH(\mathbf{p})C^\dagger &= \tau_x \kla\left( \epsilon(\mathbf{p}) \tau_0 + d_\alpha (\mathbf{p}) \tau_\alpha \right) \tau_x = \\
%&= \epsilon(\mathbf{p}) \tau_0 + d_\alpha (\mathbf{p}) \tau_x \kla\left( \tau_x\tau_\alpha + \klam\left[ \tau_\alpha, \tau_x \right] \right) = \\
%&= H(\mathbf{p}) + d_\alpha (\mathbf{p}) \tau_x 2\i\epsilon_{\alpha x \beta} \tau_\beta = \\
%&= H(\mathbf{p}) + 2\i d_\alpha (\mathbf{p}) \epsilon_{\alpha x\beta} \kla\left( \delta_{x\beta}\tau_0 + \i\epsilon_{x\beta\gamma} \tau_\gamma \right) = \\
%&=  H(\mathbf{p}) -2 d_\alpha (\mathbf{p}) 
%\end{align}
%\begin{align}
%CH(\mathbf{p})C^\dagger &= \tau_x \kla\left( \epsilon(\mathbf{p}) \tau_0 + d_\alpha (\mathbf{p}) \tau_\alpha \right) \tau_x = \epsilon(\mathbf{p}) \tau_0 + d_\alpha (\mathbf{p}) \tau_x \kla\left( \i\epsilon_{\alpha x\beta}\tau_\beta \right) = \\
%&= 
%\end{align}

\newpage
\section*{Problem 4.2: Bogoliubov-de Gennes Hamiltonians}
We will consider three Hamiltonians of differing dimensionality.
Let $\tau_\alpha$ and $\sigma_\alpha$ be Pauli-matrices for the spin and orbital degrees of freedom respectively.
Note that in general, we may write
\begin{align}
H &= \sum_{p\sigma} \epsilon_p c_{p\sigma}^\dagger c_{p\sigma} = \frac{1}{2}\sum_{p\sigma} \epsilon_p c_{p\sigma}^\dagger c_{p\sigma} + \epsilon_p \kla\left( \klammer\left\{ c_{p\sigma},c_{p\sigma}^\dagger \right\} - c_{p\sigma}c_{p\sigma}^\dagger \right) = \\
&= \frac{1}{2} \sum_{p\sigma} \epsilon_p c_{p\sigma}^\dagger c_{p\sigma} - \epsilon_p c_{p\sigma}c_{p\sigma}^\dagger + \underbrace{\frac{1}{2}\sum_{p\sigma} \epsilon_p}_{=\,\sum_p\epsilon_p \,\equiv\, \tecxt\text{const}}.
\end{align}

\subsection*{a) One-dimensional}
Consider a Bloch Hamiltonian of the form
\begin{align}
H_{1\tecxt\text{D}} &= \kla\left( p - 1 \right)\sigma_0 \equiv \epsilon_p\sigma_0,
\end{align}
then the actual Hamiltonian is given by
\begin{align}
H &= \sum_{p} c_p^\dagger  H_{1\tecxt\text{D}}(p) c_p = \frac{1}{2} \sum_p \epsilon_pc_p^\dagger c_p - \epsilon_p c_pc_p^\dagger + \tecxt\text{const},
\end{align}
where $c_p = (c_{p\uparrow}, c_{p\downarrow})^\top$. 
Including an arbitrary interaction-term of the form
\begin{align}
H_\Delta &= \sum_{p\sigma\sigma'} \Delta_{p\sigma\sigma'} c_{p\sigma}^\dagger c_{-p\sigma'}^\dagger + \tecxt\text{h.c.},
\end{align}
we can transform the Hamiltonian, also dropping the constant term.
Introducing $a_p := (c_{p\uparrow}, c_{p\downarrow}, c_{-p\uparrow}^\dagger, c_{-p\downarrow}^\dagger)^\top$ we may express the Hamiltonian as
\begin{align}
\tilde{H} &= \frac{1}{2}\sum_p a_p^\dagger H_\tecxt\text{BdG} a_p,
\end{align}
where the Bogoliubov-de\,Gennes-Hamiltonian is given by
\begin{align}
H_\tecxt\text{BdG} &= \begin{pmatrix}
H_{1\tecxt\text{D}} & \Delta_p \\ \Delta_p^\dagger & -H_{1\tecxt\text{D}}
\end{pmatrix}.
\end{align}
Inserting this we find
\begin{align}
\tilde{H} &= \frac{1}{2}\sum_p \begin{pmatrix}
c_p^\dagger & c_{-p}
\end{pmatrix} \begin{pmatrix}
\epsilon_p & \Delta_p \\ \Delta_p^\dagger & -\epsilon_p
\end{pmatrix} \begin{pmatrix}
 c_p \\ c_{-p}^\dagger
\end{pmatrix} = \\
&= \frac{1}{2}\sum_p c_p^\dagger \epsilon_p c_p - c_{-p} \epsilon_p c_{-p}^\dagger + c_p^\dagger \Delta_p c_{-p}^\dagger + \tecxt\text{h.c.} = \\
&= H + \tecxt\text{const} + \sum_{p\sigma} \Delta_{p\sigma\sigma'} c_{p\sigma}^\dagger c_{-p\sigma'}^\dagger + \tecxt\text{h.c.} \equiv H + H_\Delta.
\end{align}
It may be beneficial to choose
\begin{align}
\Delta_p &:= \begin{pmatrix}
0 & \Delta \\ -\Delta & 0
\end{pmatrix},
\end{align}
since this leads to an interaction-term of the form
\begin{align}
H_\Delta &= \sum_{p\sigma\sigma'} \Delta_{p\sigma\sigma'} c_{p\sigma}^\dagger c_{-p\sigma'}^\dagger + \tecxt\text{h.c.} = \sum_p \Delta c_{p\uparrow}^\dagger c_{-p\downarrow}^\dagger - \Delta c_{p\downarrow}^\dagger c_{-p\uparrow}^\dagger + \tecxt\text{h.c.} = \\
&= \sum_p\Delta \kla\left( c_{p\uparrow}^\dagger c_{-p\downarrow}^\dagger - c_{p\downarrow}^\dagger c_{-p\uparrow}^\dagger \right) + \tecxt\text{h.c.}
%&= \sum_p \Delta \kla\left( c_{p\uparrow}^\dagger c_{-p\downarrow}^\dagger + c_{-p\uparrow}c_{p\downarrow} \right) - \Delta \kla\left( c_{p\downarrow}^\dagger c_{-p\uparrow}^\dagger + c_{-p\downarrow}c_{p\uparrow} \right).
\end{align}
Thus we find
\begin{align}
H_\tecxt\text{BdG}  &= \begin{pmatrix}
(p-1) & 0 & 0 & \Delta \\
0 & (p-1)  & -\Delta & 0 \\
0 & -\Delta^\ast & -(p-1) & 0 \\
\Delta^\ast & 0 & 0 & -(p-1)
\end{pmatrix}
\end{align}
Note that 
\begin{align}
\sigma_x\begin{pmatrix}
a&b \\ c&d
\end{pmatrix} \sigma_x &= \sigma_x\begin{pmatrix}
b & a \\ d & c
\end{pmatrix} = \begin{pmatrix}
d & c \\ b & a.
\end{pmatrix}
\end{align}
This allows us to verify the particle-holy symmetry of the BdG-Hamiltonian
\begin{align}
C H_\tecxt\text{BdG}C &= \begin{pmatrix}
0 & \sigma_x \\ \sigma_x & 0
\end{pmatrix}\begin{pmatrix}
H_{1\tecxt\text{D}} & \Delta_p \\ \Delta_p^\dagger & -H_{1\tecxt\text{D}}
\end{pmatrix}\begin{pmatrix}
0 & \sigma_x \\ \sigma_x & 0
\end{pmatrix} = \\
&=  \begin{pmatrix}
0 & \sigma_x \\ \sigma_x & 0
\end{pmatrix} \begin{pmatrix}
\Delta_p\sigma_x & H_{1\tecxt\text{D}}\sigma_x \\ 
-H_{1\tecxt\text{D}}\sigma_x & \Delta_p^\dagger\sigma_x
\end{pmatrix} = \\
&= \begin{pmatrix}
-\sigma_x H_{1\tecxt\text{D}}\sigma_x & \sigma_x\Delta_p^\dagger \sigma_x \\ 
\sigma_x \Delta_p\sigma_x & \sigma_xH_{1\tecxt\text{D}}\sigma_x
\end{pmatrix} = \\
&= \begin{pmatrix}
-H_{1\tecxt\text{D}} & \Delta_p^\ast \\ 
-\Delta_p & H_{1\tecxt\text{D}}
\end{pmatrix} = ???
\end{align}


%\begin{thebibliography}{9}
%\bibitem{example} description, \url{https://www.exmapleurl}, day.\,month~2021.
%\end{thebibliography}

\end{document}